\section{符号说明}

\begin{table}[H]
\centering
\caption{主要集合、参数与变量}
\label{tab:symbols}
\begin{tabularx}{\textwidth}{l X}
\toprule
符号 & 含义 \\
\midrule
$e\in\Ecal$ & 患者申请事件，含来源、患者、释放时刻和截止时刻 \\
$j\in\J_e$ & 事件$e$的服务项目；$n_e=|\J_e|$ \\
$r\in\R$ & 超声设备或对应检查房间，$|\R|=22$ \\
$t\in\T$ & 5分钟离散时隙；每天上午、下午各48个 \\
$q_e\in\{I,O,P\}$ & 事件来源：住院、门诊、体检 \\
$a_e,d_e$ & 事件释放时刻与截止时刻 \\
$p_{ej}^{(0)},p_{ej}^{(1)}$ & 项目名义时长和特殊病例上界时长（时隙数） \\
$z_e$ & 特殊病例标记，$\sum_e z_e/|\Ecal|=5\%$ \\
$\R_{ej}^{S},\R_{ej}^{F}$ & 项目的严格兼容设备集合与类别回退设备集合 \\
$\R_{ej}$ & 主情景可用设备集合；床旁项目固定为$\{7\}$ \\
$c_{wt}$ & 星期$w$、日内时隙$t$的医生并发容量 \\
$b_{ej},f_{ej},u_{ej}$ & 床旁、空腹、憋尿标记 \\
$\tau$ & 相邻项目跨室转运时长，主情景为10分钟 \\
$x_{ejrt}$ & 项目$(e,j)$是否在设备$r$、时隙$t$开始 \\
$y_e$ & 事件$e$是否全部项目在截止前完成 \\
$C_e,S_e$ & 事件完成时刻与首个项目开始时刻 \\
$L_r(t)$ & 时刻$t$前设备$r$的累计占用时隙数 \\
$\kappa_r$ & 设备$r$的需求稀缺度权重 \\
\bottomrule
\end{tabularx}
\end{table}

为避免与自然常数混淆，正文用$e$表示一个具体事件，用$\Ecal$表示事件集合；集合上标$S$和$F$分别表示严格证据与回退证据。完成率均按事件计算而非项目行计算：只有$n_e$个服务项目全部排入且最后完成时刻不晚于$d_e$，才有$y_e=1$。
